\chapter{Circuits: Loops, Series, Parallel}\label{ch:g7:circuits-series-parallel}

Four years of circuits have filled your toolbox: loops, diagrams,
the two great wirings. But one character has worked namelessly all
along --- the something that flows. This year it takes its name,
its direction, and its strangest property: it is never used up.
With the current on stage, series and parallel stop being recipes
and become laws.

\section{The current takes the stage}

\begin{definition}[Electric current]\label{def:g7:circuits-series-parallel:current}
The \emph{electric current}\index{electric current} is the flow
that runs around a \omterm{def:g3:first-electric-circuit:circuit}{closed circuit}: a crowd of charged particles
--- unimaginably many, unimaginably small, cousins of the
molecular world --- set marching through the metal of the wires
by the \omterm{def:g3:first-electric-circuit:battery}{battery}'s push. No loop, no march: the current exists only
while the road is closed. (What exactly the marchers are, and how
a metal lets them through, is a fine story reserved for the last
years of this book.)
\end{definition}

\begin{proposition}[The conventional direction]\label{prop:g7:circuits-series-parallel:direction}
By worldwide convention, \omterm{def:g6:circuits-and-safety:diagram}{circuit diagrams} mark the current as
flowing \emph{outside the \omterm{def:g3:first-electric-circuit:battery}{battery} from the $+$ \omterm{def:g3:first-electric-circuit:battery}{terminal}, around
the circuit, to the $-$ \omterm{def:g3:first-electric-circuit:battery}{terminal}}. Every arrow on every diagram
in the world obeys this convention --- chosen long before anyone
could ask the marchers themselves. (When physics finally could
ask, the answer held a famous small joke, saved for a later
year.)
\end{proposition}

\begin{omfigure}
\begin{tikzpicture}[scale=0.9]
  \draw (0,0) to[battery1, l=battery] (0,3.0)
    to[short] (2.0,3.0)
    to[lamp] (4.4,3.0)
    to[short] (6.4,3.0)
    to[short] (6.4,0)
    to[short] (0,0);
  \draw[-Stealth, very thick, omThm] (1.1,3.25) -- (1.9,3.25);
  \draw[-Stealth, very thick, omThm] (6.65,1.9) -- (6.65,1.1);
  \draw[-Stealth, very thick, omThm] (5.3,-0.25) -- (4.5,-0.25);
  \draw[-Stealth, very thick, omThm] (1.9,-0.25) -- (1.1,-0.25);
  \node[right, font=\small, omThm] at (7.0,1.5)
    {the current: out of $+$, around, into $-$};
\end{tikzpicture}

{\small The convention every electrician on Earth shares: current
arrows leave the \omterm{def:g3:first-electric-circuit:battery}{battery}'s $+$, tour the loop, and return by
$-$.}
\end{omfigure}

\begin{proposition}[The current is not consumed]\label{prop:g7:circuits-series-parallel:notconsumed}
Devices do not eat the current. Whatever current enters a lamp,
\omterm{ex:g6:circuits-and-safety:motor}{motor} or \omterm{ex:g6:circuits-and-safety:motor}{buzzer} leaves it entire, marching on: what the device
takes is \emph{\omterm{def:g5:everyday-energy:energy}{energy}} from the \omterm{def:g3:first-electric-circuit:battery}{battery}'s store, delivered by the
passing march. In any series loop, one and the same current
therefore flows through every device --- as much arrives back at
the \omterm{def:g3:first-electric-circuit:battery}{battery} as ever left it.
\end{proposition}

\begin{example}[The misconception on trial]\label{ex:g7:circuits-series-parallel:trial}
The tempting error: ``the first \omterm{def:g3:first-electric-circuit:bulb}{bulb} uses up some current, so the
second must glow weaker.'' The circuit itself refutes it. Wire
two \emph{identical} \omterm{def:g3:first-electric-circuit:bulb}{bulbs} \omterm{def:g4:series-circuits:series}{in series} and look: both glow with
exactly equal brightness --- and swapping them changes nothing.
If current were consumed on the way, \omterm{def:g3:first-electric-circuit:bulb}{bulb} two should always be
the dim one, whichever \omterm{def:g3:first-electric-circuit:bulb}{bulb} sits there. It never is. The current
arrives whole; what the pair shares is the \omterm{def:g3:first-electric-circuit:battery}{battery}'s \emph{push},
which is why both are dimmer than a lone \omterm{def:g3:first-electric-circuit:bulb}{bulb} --- equally dimmer.
\end{example}

\section{The two families, as laws}

\begin{proposition}[The series law]\label{prop:g7:circuits-series-parallel:series}
In a \omterm{def:g4:series-circuits:series}{series circuit} --- one single loop --- the same current
flows through every device: one road, one march, no exceptions.
Consequences, now explained: one gap stops the march everywhere
at once; identical devices share the work equally, wherever they
sit; and the order of devices on the loop cannot matter to a
march that passes through all of them alike.
\end{proposition}

\begin{proposition}[The parallel law]\label{prop:g7:circuits-series-parallel:parallel}
At a \omterm{def:g5:series-parallel-circuits:parallel}{junction}, the current \emph{divides}: part of the march
takes each \omterm{def:g5:series-parallel-circuits:parallel}{branch}, and the parts rejoin, complete, at the far
\omterm{def:g5:series-parallel-circuits:parallel}{junction} --- nothing lost, nothing gained. Each \omterm{def:g5:series-parallel-circuits:parallel}{branch} receives
the \omterm{def:g3:first-electric-circuit:battery}{battery}'s full push, so each device works as if alone; and a
blocked \omterm{def:g5:series-parallel-circuits:parallel}{branch} merely sends nobody, while the other \omterm{def:g5:series-parallel-circuits:parallel}{branches}
march on.
\end{proposition}

\begin{omfigure}
\begin{tikzpicture}[scale=0.85]
  \draw (0,0) to[battery1] (0,3.2)
    to[short] (2.4,3.2);
  \draw (2.4,3.2) to[short] (2.4,2.3)
    to[lamp] (5.6,2.3) to[short] (5.6,3.2);
  \draw (2.4,3.2) to[short] (2.4,4.1)
    to[lamp] (5.6,4.1) to[short] (5.6,3.2);
  \draw (5.6,3.2) to[short] (7.2,3.2) to[short] (7.2,0)
    to[short] (0,0);
  \fill (2.4,3.2) circle (2pt);
  \fill (5.6,3.2) circle (2pt);
  \draw[-Stealth, very thick, omThm] (1.0,3.45) -- (1.9,3.45);
  \draw[-Stealth, thick, omThm] (2.85,4.35) -- (3.65,4.35);
  \draw[-Stealth, thick, omThm] (2.85,2.05) -- (3.65,2.05);
  \draw[-Stealth, very thick, omThm] (6.05,3.45) -- (6.95,3.45);
  \node[above, font=\small, omThm] at (1.45,3.55) {whole march};
  \node[above, font=\small, omThm] at (6.5,3.55) {whole again};
\end{tikzpicture}

{\small The parallel law drawn: the march divides at the first
\omterm{def:g5:series-parallel-circuits:parallel}{junction}, each part tours its \omterm{def:g5:series-parallel-circuits:parallel}{branch}, and the full march
reassembles at the second.}
\end{omfigure}

\begin{example}[Reading brightness like an electrician]\label{ex:g7:circuits-series-parallel:brightness}
One \omterm{def:g3:first-electric-circuit:battery}{battery}, identical \omterm{def:g3:first-electric-circuit:bulb}{bulbs}. Alone: full brightness --- the
reference. Two \omterm{def:g4:series-circuits:series}{in series}: both equally dim --- one march, shared
push. Two \omterm{def:g5:series-parallel-circuits:parallel}{in parallel}: both fully bright --- each \omterm{def:g5:series-parallel-circuits:parallel}{branch} enjoys
the whole push (and the \omterm{def:g3:first-electric-circuit:battery}{battery} empties roughly twice as fast:
nothing is free). Three \omterm{def:g4:series-circuits:series}{in series}: dimmer still, all equal. The
electrician's glance: count the loop-mates for dimness, count the
\omterm{def:g5:series-parallel-circuits:parallel}{branches} for \omterm{def:g3:first-electric-circuit:battery}{battery} bills.
\end{example}

\begin{method}[Predicting any lamp's fate]\label{met:g7:circuits-series-parallel:predict}
Given a diagram, for each lamp:
\begin{enumerate}
  \item \emph{Existence}: is there at least one unbroken loop
        from $+$ through this lamp to $-$? No loop, no \omterm{def:g1:light-and-shadows:source}{light}.
  \item \emph{Company on the road}: which devices share
        \emph{every} loop through this lamp? Those are its series
        companions --- more companions, dimmer glow.
  \item \emph{\omterm{def:g5:series-parallel-circuits:parallel}{Branch} neighbors}: which devices sit on other
        \omterm{def:g5:series-parallel-circuits:parallel}{branches}? They neither dim this lamp nor darken with it.
  \item \emph{\omterm{ex:g3:first-electric-circuit:switch}{Switches}}: an open \omterm{ex:g3:first-electric-circuit:switch}{switch} kills exactly the loops
        it sits on --- re-run step 1 with it removed.
\end{enumerate}
\end{method}

\begin{remark}[Why the misconception dies hard]\label{rem:g7:circuits-series-parallel:whyhard}
``Something must get used up --- the \omterm{def:g3:first-electric-circuit:battery}{battery} empties!'' True: the
\omterm{def:g3:first-electric-circuit:battery}{battery}'s stored \emph{\omterm{def:g5:everyday-energy:energy}{energy}} is spent, handed to lamps as \omterm{def:g1:light-and-shadows:source}{light}
and warmth by the passing march. But the marchers themselves
circulate unharmed, like ski-lift chairs delivering skiers
uphill all day without a chair ever being consumed. Keep the two
bookkeepings apart --- current circulates, \omterm{def:g5:everyday-energy:energy}{energy} is delivered
--- and every circuit in this book stays honest. The instrument
that counts the march itself joins us next year.
\end{remark}

\section{Exercises}

\begin{exercise}[$\star$]\label{exo:g7:circuits-series-parallel:1}
What is the \omterm{def:g7:circuits-series-parallel:current}{electric current}? When does it exist in a circuit,
and when not?
\end{exercise}

\begin{exercise}[$\star$]\label{exo:g7:circuits-series-parallel:2}
State the conventional direction of the current. On a diagram of
a \omterm{def:g3:first-electric-circuit:bulb}{one-bulb} loop, where do the arrows point?
\end{exercise}

\begin{exercise}[$\star$]\label{exo:g7:circuits-series-parallel:3}
Two identical \omterm{def:g3:first-electric-circuit:bulb}{bulbs} \omterm{def:g4:series-circuits:series}{in series} glow exactly equally. What popular
error does this equality refute, and how?
\end{exercise}

\begin{exercise}[$\star$]\label{exo:g7:circuits-series-parallel:4}
In a series loop with a \omterm{ex:g6:circuits-and-safety:motor}{motor} and a lamp, the current through
the \omterm{ex:g6:circuits-and-safety:motor}{motor} is compared with the current through the lamp. What
does the series law say?
\end{exercise}

\begin{exercise}[$\star$]\label{exo:g7:circuits-series-parallel:5}
What happens to the current at a \omterm{def:g5:series-parallel-circuits:parallel}{parallel circuit}'s first
\omterm{def:g5:series-parallel-circuits:parallel}{junction}? And at the second? What is true of each \omterm{def:g5:series-parallel-circuits:parallel}{branch}'s share
of the \omterm{def:g3:first-electric-circuit:battery}{battery}'s push?
\end{exercise}

\begin{exercise}[$\star$]\label{exo:g7:circuits-series-parallel:6}
One \omterm{def:g3:first-electric-circuit:battery}{battery}, identical \omterm{def:g3:first-electric-circuit:bulb}{bulbs}: rank the brightness of a \omterm{def:g3:first-electric-circuit:bulb}{bulb} ---
alone; \ with one series companion; \ with one parallel
neighbor; \ with two series companions.
\end{exercise}

\begin{exercise}[$\star$]\label{exo:g7:circuits-series-parallel:7}
Which arrangement empties the \omterm{def:g3:first-electric-circuit:battery}{battery} faster: two \omterm{def:g3:first-electric-circuit:bulb}{bulbs} \omterm{def:g4:series-circuits:series}{in
series}, or the same two \omterm{def:g5:series-parallel-circuits:parallel}{in parallel}? Why, in march-language?
\end{exercise}

\begin{exercise}[$\star$]\label{exo:g7:circuits-series-parallel:8}
Explain the ski-lift picture of
\cref{rem:g7:circuits-series-parallel:whyhard}: what plays the
chairs, what plays the skiers, and what is truly used up?
\end{exercise}

\begin{exercise}[$\star\star$]\label{exo:g7:circuits-series-parallel:9}
A diagram: \omterm{def:g3:first-electric-circuit:battery}{battery}, \omterm{ex:g3:first-electric-circuit:switch}{switch}, then lamp A; after A a \omterm{def:g5:series-parallel-circuits:parallel}{junction}
sends two \omterm{def:g5:series-parallel-circuits:parallel}{branches} --- lamp B on one, lamp C on the other ---
rejoining before returning to the \omterm{def:g3:first-electric-circuit:battery}{battery}. Using
\cref{met:g7:circuits-series-parallel:predict}: who are A's
series companions? What does A's glow do when C's \omterm{def:g5:series-parallel-circuits:parallel}{branch} is cut?
\end{exercise}

\begin{exercise}[$\star\star$]\label{exo:g7:circuits-series-parallel:10}
Decorative \omterm{def:g1:light-and-shadows:source}{lights} ask for a design: ten \omterm{def:g3:first-electric-circuit:bulb}{bulbs} that all glow at
full, independent brightness from one supply --- but the whole
set must obey a single master \omterm{ex:g3:first-electric-circuit:switch}{switch}. Draw or describe the
wiring, citing both laws.
\end{exercise}

\begin{exercise}[$\star\star$]\label{exo:g7:circuits-series-parallel:11}
In the \omterm{def:g3:first-electric-circuit:bulb}{two-bulb} \omterm{def:g5:series-parallel-circuits:parallel}{parallel circuit}, \omterm{def:g3:first-electric-circuit:bulb}{bulb} 1's \omterm{def:g5:series-parallel-circuits:parallel}{branch} is thick
copper wire while \omterm{def:g3:first-electric-circuit:bulb}{bulb} 2's \omterm{def:g5:series-parallel-circuits:parallel}{branch} includes a long, thin wire.
\omterm{def:g3:first-electric-circuit:bulb}{Bulb} 2 glows a little weaker. What does this suggest about how a
\omterm{def:g5:series-parallel-circuits:parallel}{branch}'s road itself can throttle its share of the march? (Next
year gives this throttling a name and a law.)
\end{exercise}

\begin{exercise}[$\star\star\star$]\label{exo:g7:circuits-series-parallel:12}
A \omterm{def:g3:first-electric-circuit:battery}{battery} feeds two parallel \omterm{def:g5:series-parallel-circuits:parallel}{branches}: \omterm{def:g5:series-parallel-circuits:parallel}{branch} one holds lamp X
alone; \omterm{def:g5:series-parallel-circuits:parallel}{branch} two holds lamps Y and Z \omterm{def:g4:series-circuits:series}{in series}. Compare the
three lamps' brightnesses, in order, with reasons from both
laws. Then predict every lamp's fate if Z burns out --- and if
instead X burns out.
\end{exercise}

\section{Problem: The Lighthouse Keeper's Panel}

\begin{problem}[{Weekend problem --- rewiring the old lighthouse;
one battery bank, many duties; the keeper's night of
faults}]\label{pb:g7:circuits-series-parallel:1}
The lighthouse's \omterm{def:g3:first-electric-circuit:battery}{battery} bank must feed: the great lamp, a
foghorn (a mighty \omterm{ex:g6:circuits-and-safety:motor}{buzzer}), the keeper's reading \omterm{def:g1:light-and-shadows:source}{light}, and a
staircase \omterm{def:g1:light-and-shadows:source}{light} --- and survive the keeper's troubleshooting
night.

\medskip\noindent\textbf{Part I --- Design.}
\begin{enumerate}
  \item The great lamp must never dim because something else
        switches on. What wiring does this demand for the whole
        installation, and why would any series pairing betray
        the ships?
  \item Each duty needs its own \omterm{ex:g3:first-electric-circuit:switch}{switch}, and the keeper wants one
        master cut-off for storms. Place all five \omterm{ex:g3:first-electric-circuit:switch}{switches}.
  \item The foghorn and the staircase \omterm{def:g1:light-and-shadows:source}{light} must never be
        prevented from working by each other's faults. Is this
        already guaranteed? By which law?
  \item Sketch or describe the full diagram: bank, master
        \omterm{ex:g3:first-electric-circuit:switch}{switch}, four \omterm{def:g5:series-parallel-circuits:parallel}{branches}, four \omterm{def:g5:series-parallel-circuits:parallel}{branch} \omterm{ex:g3:first-electric-circuit:switch}{switches}.
\end{enumerate}

\medskip\noindent\textbf{Part II --- The keeper's questions.}
\begin{enumerate}[resume]
  \item The keeper worries: ``Four duties drinking at once ---
        does each \omterm{def:g5:series-parallel-circuits:parallel}{branch} still get the bank's full push?''
        Answer with the parallel law.
  \item ``And is current being used up as it works through the
        great lamp?'' Set the keeper straight, with the ski-lift
        if it helps.
  \item On a foggy night everything runs at once. What is the
        cost of parallel generosity, and where in the lighthouse
        is it paid?
  \item The reading \omterm{def:g1:light-and-shadows:source}{light} alone is lit; measure-minded, the
        keeper follows the march: describe its full round trip,
        naming every component it passes and every one it
        avoids.
\end{enumerate}

\medskip\noindent\textbf{Part III --- The night of faults.}
\begin{enumerate}[resume]
  \item Midnight: the staircase \omterm{def:g3:first-electric-circuit:bulb}{bulb} burns out. What else goes
        dark? Which law answers?
  \item One o'clock: a clumsy repair leaves a stretch of bare
        wire joining the two \omterm{def:g3:first-electric-circuit:battery}{terminals} of the bank directly,
        upstream of every \omterm{def:g5:series-parallel-circuits:parallel}{branch}. Name the event, and explain
        why every duty in the lighthouse now fails at once ---
        where does the march go?
  \item Two o'clock, fault cleared: the keeper wants the great
        lamp protected so that such a shortcut can never again
        silence it. A second, independent \omterm{def:g3:first-electric-circuit:battery}{battery} bank is
        available. Propose the wiring.
  \item Dawn: write the keeper's log --- three sentences: which
        law kept ships safe tonight, which misconception the
        night refuted, and which single wire caused the worst
        hour.
\end{enumerate}
\end{problem}
